\NormalFont\ShortTitle{Iuda}

{\MT Scrisoarea lui Iuda


\par }\OneChap {\PP \VerseOne{1}Iuda, rob al lui Isus Hristos și frate al lui Iacov, către cei chemați, sfințițiți de Dumnezeu Tatăl și păstrați pentru Isus Hristos:

\VS{2}Fie ca mila, pacea și dragostea să vă fie înmulțite.


\par }{\PP \VS{3}Preaiubiților, deși eram foarte nerăbdător să vă scriu despre mântuirea noastră comună, am fost silit să vă scriu pentru a vă îndemna să vă luptați cu tărie pentru credința care a fost dată sfinților o dată pentru totdeauna.

\VS{4}Pentru că sunt unii oameni care s-au strecurat pe ascuns, chiar cei despre care s-a scris demult pentru această condamnare: oameni nelegiuiți, care transformă harul Dumnezeului nostru în indecență și neagă pe singurul nostru Stăpân, Dumnezeu și Domn, Isus Hristos.


\par }{\PP \VS{5}Vreau să vă amintesc, deși știți deja, că Domnul, după ce a izbăvit un popor din țara Egiptului, a nimicit pe cei ce nu au crezut.

\VS{6}Pe îngerii care nu și-au păstrat primul domeniu, ci și-au părăsit propria locuință, i-a ținut în legături veșnice, sub întuneric, pentru judecata din ziua cea mare.

\VS{7}La fel ca Sodoma și Gomora și orașele din jurul lor, care, la fel ca acestea, s-au dedat la imoralitate sexuală și au umblat după trupuri străine, sunt arătate ca exemplu, suferind pedeapsa focului veșnic.

\VS{8}Și totuși, în același mod, și aceștia, în visteria lor, spurcă carnea, disprețuiesc autoritatea și calomniază ființele cerești.

\VS{9}Dar Mihail, arhanghelul Mihail, când se certa cu diavolul și discuta despre trupul lui Moise, nu a îndrăznit să aducă împotriva lui o condamnare abuzivă, ci a spus: “Domnul să te certe!”.

\VS{10}Aceștia însă vorbesc de rău despre orice lucru pe care nu-l cunosc. Ei sunt distruși în aceste lucruri pe care le înțeleg în mod natural, asemenea creaturilor fără rațiune.

\VS{11}Vai de ei! Pentru că au mers pe calea lui Cain, au alergat cu destrăbălare în rătăcirea lui Balaam, pentru a fi plătiți, și au pierit în răzvrătirea lui Core.

\VS{12}Aceștia sunt recife stâncoase ascunse în ospețele voastre de dragoste, când se ospătează cu voi, păstori care fără teamă se hrănesc singuri; nori fără apă, purtați de vânturi; pomi de toamnă fără rod, morți de două ori, smulși din rădăcini;

\VS{13}valuri sălbatice ale mării, care își scot spuma rușinii lor; stele rătăcitoare, cărora le-a fost rezervată pentru totdeauna negura întunericului.

\VS{14}Despre acestea a profețit și Enoh, al șaptelea de la Adam, care a spus: “Iată, Domnul a venit cu zece mii de sfinți ai săi,

\VS{15}pentru a judeca pe toți și pentru a-i condamna pe toți cei nelegiuiți de toate faptele lor de impietate pe care le-au făcut în mod nelegiuit și de toate lucrurile grele pe care păcătoșii nelegiuiți le-au spus împotriva lui”.

\VS{16}Aceștia sunt murmurători și plângăcioși, umblând după poftele lor — și gura lor vorbește lucruri mândre — arătând respect de persoane pentru a obține avantaje.


\par }{\PP \VS{17}Dar voi, preaiubiților, aduceți-vă aminte de cuvintele care au fost rostite mai înainte de apostolii Domnului nostru Isus Hristos.

\VS{18}Ei v-au spus: “În vremurile din urmă vor fi batjocoritori, care vor umbla după poftele lor nelegiuite.”

\VS{19}Aceștia sunt cei care provoacă dezbinare și sunt senzuali, neavând Duhul Sfânt.


\par }{\PP \VS{20}Iar voi, iubiților, zidiți-vă pe credința voastră cea sfântă, rugându-vă în Duhul Sfânt.

\VS{21}Păstrați-vă în dragostea lui Dumnezeu, așteptând îndurarea Domnului nostru Isus Hristos pentru viața veșnică.

\VS{22}Pe unii îi compătimește, făcând distincție,

\VS{23}iar pe alții îi mântuiește, smulgându-i cu frică din foc, urând chiar și hainele pătate de carne.


\par }{\PP \VS{24}Iar Celui ce poate să-i păzească de poticnire și să vă înfățișeze fără greșeală în fața slavei Lui, cu mare bucurie,

\VS{25}lui Dumnezeu, Mântuitorul nostru, singurul înțelept, fie slava și maiestatea, stăpânirea și puterea, acum și pururea și în vecii vecilor. Amin.


\par }